\documentclass[11pt, a4paper, ngerman]{article}
\usepackage[utf8]{inputenc}
\usepackage{babel}
\usepackage{amsmath}
\usepackage{hyperref}
\usepackage{geometry}
\geometry{a4paper, margin=2.5cm}

\title{Review: Photosynthesis -- Mechanisms, Efficiency, and Relevance for Renewable Energy Systems (Paper von Eliah Graf)}
\author{Luka Tadic}
\date{\today}

\begin{document}
\maketitle

\section{Zusammenfassung}
Das Paper behandelt die Photosynthese aus einer ingenieurwissenschaftlichen Perspektive mit Fokus auf dem Vergleich zur Photovoltaik und künstlicher Photosynthese. Der rote Faden ist logisch und gut nachvollziehbar aufgebaut.

\section{Formelle Kriterien und Spezifikationen}
\begin{itemize}
    \item \textbf{Format:} Das IEEE-Template wurde ordnungsgemäß angewendet.
    \item \textbf{Länge:} Das Paper umfasst 4 Seiten, da das Literaturverzeichnis (die Quellen) zur Gesamtlänge dazuzählt. Die formelle Längenvorgabe von 4--6 Seiten ist somit erfüllt.
    \item \textbf{Abbildungen:} Es sind 3 Abbildungen und eine Tabelle inkludiert. Die geforderte Mindestzahl von 4--6 integrativ eingebundenen Abbildungen (inklusive Tabellen) ist somit nun erfüllt.
\end{itemize}

\section{Inhaltliche und Technische Tiefe (,,Storytelling`)}
\textbf{Storytelling:} Das Storytelling ist hervorragend. Der Bogen von biologischen Grundlagen (Z-Scheme) über die technische Evaluierung (Efficiency Analysis) hin zu Zukunftstechnologien (Artificial Photosynthesis) ist methodisch sinnvoll und klar strukturiert. \\\\
\textbf{Technische Tiefe:} Der geforderte Fokus auf die Technik und die ingenieurwissenschaftliche Durchdringung wurden sehr gut umgesetzt, vor allem im Vergleich zu PV-Zellen. Inhaltlich ist die technische Tiefe vollkommen überzeugend und das passende Verhältnis von technischem Inhalt zu den biologischen Grundlagen ist hervorragend abgedeckt.

\section{Quellen und Zitationen}
\begin{itemize}
    \item \textbf{Quantität:} Die Anzahl der Quellen erfüllt mit $>$ 9 Quellen (\textit{hier: 22 Quellen}) die Anforderungen.
    \item \textbf{Qualität:} Die Qualität einiger Quellen (z.\,B. \url{wikipedia}, Mitschriften wie \textit{BioNinja} etc.) muss kritisiert werden. In einem wissenschaftlichen Kontext sollten primär Paper oder Lehrbücher zitiert werden, besonders im Hinblick auf Halluzinationsfreiheit und Verifizierbarkeit. Die Überprüfung und Bereinigung der Quellen ist der wichtigste handwerkliche Punkt für eine letzte Überarbeitung.
    \item \textbf{Erreichbarkeit (DOIs und Links):} Viele Hyperlinks der verwendeten Online-Quellen (wie z.\,B. von libretexts.org) leiten ins Leere oder sind fehlerhaft. Es ist zwingend darauf zu achten, dass \textbf{DOIs} bei den Quellen vorhanden sind. Dies stellt sicher, dass Leser jederzeit und dauerhaft auf das zugrunde liegende Material zugreifen können.
    \item \textbf{Interne Verlinkung (\textit{Best Practice}):} Es ist eine sehr gute Praxis, Zitationsnummern im Text direkt mit dem Literaturverzeichnis zu verlinken. Wenn der Leser darauf klickt, gelangt er sofort zur jeweiligen Quelle. Da das Paper in LaTeX geschrieben wurde, bietet sich zur einfachen Umsetzung das Paket \texttt{hyperref} an.
    \item \textbf{Formatierung:} Die Formatierung der Literaturverweise nach dem IEEE-Zitierstil (z.\,B. [1] Autor, ,,Titel``, \dots) ist zulässig und in Ordnung. Eine Anpassung auf einen anderen Reference Style ist nicht notwendig; der Fokus sollte stattdessen auf der Qualität der Quellen und der Angabe von DOIs liegen.
\end{itemize}

\section{Grafiken und Schreibstil}
\begin{itemize}
    \item \textbf{Grafiken:} Die eigens erstellten Grafiken sind ein großer Pluspunkt des Papers. Sie veranschaulichen die Sachverhalte sehr gut, sind hochauflösend und farblich gut abgegrenzt. \textbf{Ein kleiner Kritikpunkt ist jedoch die erste Abbildung:} Diese ragt etwas zu sehr an bzw. über die Seitenränder hinaus. Dies sollte im Layout noch korrigiert werden.
    \item \textbf{Sprache \& Text:} Der Text ist gut lesbar, weist einen flüssigen Schreibstil auf und lässt keine offensichtlichen Wiederholungen erkennen.
\end{itemize}

\section{Fazit und konstruktive Kritik}
Das Paper ist inhaltlich stark, punktet mit hervorragenden Diagrammen und weist ein gutes Storytelling auf. Die technische Tiefe und die formale Seitenzahl sind bereits auf einem absolut passenden Stand. 

Der einzige, aber wichtige Kritikpunkt liegt bei den \textbf{Quellen} sowie einer Grafik:
\begin{enumerate}
    \item Die wissenschaftliche Qualität einiger Internetquellen sowie defekte Hyperlinks müssen dringend überprüft werden.
    \item Jedem wissenschaftlichen Paper sollte nach Möglichkeit eine DOI hinzugefügt werden, um den dauerhaften Zugriff für Leser zu garantieren.
    \item Der Seitenrand bei der ersten Abbildung sollte noch angepasst werden.
\end{enumerate}
Es wird empfohlen, diese letzte inhaltliche Überprüfung des Literaturverzeichnisses vorzunehmen, um die Anforderungen des Seminars vollumfänglich und auf einem sehr hohen wissenschaftlichen Niveau zu erfüllen.

\end{document}